\chapter{Estado da Arte}
\label{cap:estado-arte}

Este capítulo discute os trabalhos que se aproximam do problema formulado no
Capítulo~\ref{cap:proposta}. O critério de seleção foi a existência de alguma
tentativa de caracterizar o estado de generalização de uma rede sem depender
integralmente de um conjunto de teste rotulado, ou de caracterizar a transição
entre aprendizado e memorização. A Seção~\ref{sec:E-memorizacao} revisa o que se
sabe sobre memorização em redes profundas. A Seção~\ref{sec:E-pesos} discute
métricas de generalização calculadas sobre os pesos. A
Seção~\ref{sec:E-complexidade} trata da aplicação de medidas de complexidade a
redes neurais. A Seção~\ref{sec:E-parada} revisa critérios de parada. A
Seção~\ref{sec:E-sintese} sintetiza e posiciona a proposta.

\section{Memorização e Generalização em Redes Profundas}
\label{sec:E-memorizacao}

\citeonline{zhang2017} demonstraram experimentalmente que redes convolucionais
de grande porte ajustam perfeitamente rótulos atribuídos ao acaso, e que esse
comportamento persiste mesmo sob regularização explícita. O resultado é
qualitativamente importante para este trabalho por duas razões: estabelece que a
capacidade de memorização não é um acidente, mas uma propriedade das
arquiteturas usuais, e fornece o protocolo experimental --- corrupção
controlada de rótulos --- que permite induzir memorização de forma reprodutível.
A limitação, do ponto de vista aqui adotado, é que o diagnóstico continua
dependendo da comparação entre desempenho de treino e de teste.

\citeonline{arpit2017} refinam esse quadro ao mostrar que memorização e
aprendizado de padrões não ocorrem simultaneamente: redes tendem a aprender
primeiro as regularidades simples e só depois memorizam os exemplos ruidosos.
A existência de dois regimes sucessivos, e não de um processo homogêneo, é o
pressuposto que torna plausível a hipótese deste projeto --- se há transição,
pode haver assinatura da transição. Os autores, contudo, caracterizam os
regimes por meio de métricas que dependem dos dados, e não da estrutura dos
pesos.

\citeonline{power2022} descrevem o fenômeno de \emph{grokking}, em que a
generalização ocorre muito depois do ajuste completo do treino, às vezes após
longo período de estagnação da validação. O fenômeno é relevante aqui porque
constitui um contraexemplo aos critérios ingênuos de parada antecipada: uma
interrupção baseada apenas na estagnação da validação eliminaria justamente o
regime em que a generalização acabaria emergindo. Um indicador que leia o
estado interno da rede, e não apenas seu desempenho externo, é potencialmente
capaz de distinguir estagnação de reorganização --- questão que este trabalho
deixa registrada como hipótese derivada, ainda não testada.

\section{Métricas de Generalização a partir dos Pesos}
\label{sec:E-pesos}

\citeonline{bartlett2017} propõem limites de generalização baseados em normas
espectrais e margens, estabelecendo que propriedades geométricas dos pesos se
relacionam com a capacidade de generalizar. \citeonline{dziugaite2017} obtêm,
por via PAC-Bayesiana, limites não vacuosos para redes sobreparametrizadas.
Ambos representam avanços substantivos, mas produzem limites válidos para um
estado da rede, tipicamente o final do treinamento, e não uma leitura
acompanhável ao longo das épocas.

\citeonline{jiang2020} conduzem o estudo empírico mais sistemático dessa
família, avaliando um grande conjunto de medidas candidatas quanto à sua
correlação causal com a lacuna de generalização. A conclusão é cautelar: muitas
medidas teoricamente motivadas correlacionam-se mal com a generalização
observada, e medidas baseadas em otimização desempenham-se melhor do que as
baseadas em norma. Esse resultado impõe a este trabalho uma exigência
metodológica clara --- não basta exibir correlação entre indicador e
\emph{overfitting} em uma execução; é preciso demonstrar discriminação entre
execuções com e sem memorização, com controles pareados.

\citeonline{martin2021} mostram que propriedades de cauda pesada do espectro das
matrizes de pesos permitem prever tendências de qualidade de modelos sem acesso
aos dados de treino ou de teste. Este é o trabalho mais próximo do objetivo aqui
perseguido, e convém explicitar a diferença: a abordagem é \emph{espectral},
apoiando-se nos autovalores das matrizes, ao passo que a abordagem deste projeto
é \emph{informação-teórica}, apoiando-se na distribuição de valores e na
organização espacial dos pesos. Além disso, o alvo em \citeonline{martin2021} é
comparar modelos já treinados, enquanto aqui o alvo é acompanhar a trajetória
temporal de um único treinamento.

\section{Medidas de Complexidade Aplicadas a Redes Neurais}
\label{sec:E-complexidade}

A aplicação de medidas de entropia e complexidade ao estudo de redes neurais é
mais recente e menos consolidada. \citeonline{advani2020} analisam a dinâmica do
aprendizado em regime de alta dimensionalidade, caracterizando fases distintas
do treinamento. \citeonline{delgadobonal2019} revisam o uso de medidas de
entropia em séries temporais e sistematizam a escolha de parâmetros, questão
diretamente relevante para a estabilidade das estimativas aqui empregadas.

O que a revisão da literatura não localizou foi trabalho que acompanhe, época a
época, medidas de complexidade estatística calculadas sobre os pesos de uma
camada específica com o propósito explícito de emitir um aviso antecipado de
\emph{overfitting}, com controles pareados e quantificação da antecedência.
Essa é a lacuna que o projeto endereça.

\section{Critérios de Parada}
\label{sec:E-parada}

\citeonline{prechelt1998} sistematizou os critérios automáticos de parada
antecipada baseados em validação cruzada, avaliando empiricamente catorze
critérios quanto a eficiência e eficácia. O trabalho permanece a referência
metodológica da prática corrente e explicita seu próprio pressuposto: a decisão
de parada é uma função do erro medido em dados separados. É exatamente esse
pressuposto que este projeto procura relaxar, não substituindo o critério
clássico, mas oferecendo uma segunda leitura obtida por caminho independente.

\section{Síntese e Posicionamento da Proposta}
\label{sec:E-sintese}

\begin{table}[H]
    \centering
    \caption{Comparação dos trabalhos discutidos quanto à dependência de dados
    rotulados e à natureza temporal da leitura.}
    \label{tab:estado-arte}
    \footnotesize
    \begin{tabular}{lcccc}
        \toprule
        Trabalho & Objeto da medida & Precisa de rótulo & Leitura temporal & Avisa antes? \\
        \midrule
        \citeonline{prechelt1998} & Erro de validação   & Sim  & Sim & Não \\
        \citeonline{zhang2017}    & Desempenho          & Sim  & Não & Não \\
        \citeonline{arpit2017}    & Desempenho          & Sim  & Sim & Não \\
        \citeonline{bartlett2017} & Norma espectral     & Não  & Não & Não \\
        \citeonline{dziugaite2017}& Limite PAC-Bayes    & Sim  & Não & Não \\
        \citeonline{jiang2020}    & Diversas            & Sim  & Não & Não \\
        \citeonline{martin2021}   & Espectro dos pesos  & Não  & Não & Não \\
        \textbf{Este trabalho}    & Complexidade dos pesos & Não & Sim & Sim \\
        \bottomrule
    \end{tabular}
    \fonte{Elaborado pelo autor.}
\end{table}

A Tabela~\ref{tab:estado-arte} organiza os trabalhos segundo dois eixos. O
primeiro é a dependência de dados rotulados: apenas \citeonline{bartlett2017} e
\citeonline{martin2021} dispensam-nos, e ambos o fazem por via espectral. O
segundo é a natureza temporal da leitura: das abordagens que dispensam rótulos,
nenhuma foi concebida para acompanhar a trajetória do treinamento época a
época.

A proposta deste trabalho situa-se na interseção ainda vazia desses dois eixos,
e diferencia-se da literatura existente em três aspectos: (i) a medida é
informação-teórica, incidindo sobre a distribuição de valores e sobre a
organização espacial dos pesos, e não sobre seu espectro; (ii) a leitura é
temporal, produzindo uma trajetória e não um escalar ao final do treino; e
(iii) a validação é feita por controles pareados com indução explícita de
memorização, em resposta direta à advertência metodológica de
\citeonline{jiang2020} quanto a correlações espúrias.
